% Self-contained NeurIPS-style short report: no external style, bibliography,
% figures, or generated inputs. Compile with: pdflatex neurips.tex (twice).
\documentclass[10pt,letterpaper]{article}
\usepackage[T1]{fontenc}
\usepackage[utf8]{inputenc}
\usepackage{mathptmx}
\usepackage{courier}
\usepackage[textwidth=5.5in,textheight=9in,top=1in,headheight=12pt,headsep=18pt]{geometry}
\usepackage{amsmath,amssymb,booktabs,microtype}
\usepackage[colorlinks=true,linkcolor=black,citecolor=black,urlcolor=blue,
pdftitle={Replicating Intel's Acoustic Weld-Defect Detection Baseline at 16 kHz},
pdfauthor={Jaden Majid, Hieu Le, Chris Tardy, Irene Berezin}]{hyperref}
\urlstyle{same}
\setlength{\parindent}{0pt}
\setlength{\parskip}{5.5pt}
\setlength{\textfloatsep}{10pt}
\setlength{\abovecaptionskip}{5pt}
\setlength{\belowcaptionskip}{0pt}
\makeatletter
\renewcommand\section{\@startsection{section}{1}{\z@}{-2.0ex plus -.5ex}{.8ex}{\large\bfseries}}
\renewcommand\subsection{\@startsection{subsection}{2}{\z@}{-1.5ex plus -.4ex}{.5ex}{\normalsize\bfseries}}
\makeatother

\begin{document}
\begin{center}
\hrule height 4pt
\vspace{9pt}
{\LARGE\bfseries Replicating Intel's Acoustic Weld-Defect\par
Detection Baseline at 16 kHz\par}
\vspace{9pt}
\hrule height 1pt
\vspace{12pt}
{\normalsize Jaden Majid \quad Hieu Le \quad Chris Tardy \quad Irene Berezin}\\[3pt]
Grad Student Descent \quad $\cdot$ \quad AI 360\\[2pt]
{\small Experiment snapshot: November 4, 2026}
\end{center}
\vspace{-2pt}
\begin{center}\textbf{Abstract}\end{center}
\begin{quotation}\small
We reimplement Intel's one-dimensional convolutional autoencoder for acoustic
weld-defect detection using the released mono 16\,kHz recordings. The experiment
retains the published best bottleneck and approximately preserves its spectral
analysis duration, while making unspecified implementation choices explicit.
The designated primary run achieves test ROC-AUC 0.8178, compared with the
published 192\,kHz result of 0.8460. Across three completed 50-epoch runs,
mean test ROC-AUC is $0.8180\pm0.0007$ and mean validation ROC-AUC is
$0.8158\pm0.0005$ (sample standard deviations). These results establish a baseline for the released audio,
but do not isolate the effect of sampling rate or establish exact reproduction
of the original benchmark.
\end{quotation}

\section{Scope and data}
Stemmer et al.~\cite{stemmer2024} detect defective welds by reconstructing
normal welding audio. We reproduce their audio-only model from scratch as a
baseline for subsequent representation-learning experiments. This is a single
configuration with three planned random seeds, rather than a repeat of the
paper's FFT/bottleneck grid search. The configuration is chosen from the paper,
before inspecting replication test scores.

The audited release contains 4,040 mono FLAC recordings at 16\,kHz, spanning
normal welds and eleven defect categories. We retain the manifest partitions
(Table~\ref{tab:data}); only normal training recordings enter optimization.
Recordings last 25--38\,s. Validation contains labeled defects and is used to
select epochs, so training is normal-only but selection uses anomaly labels.
The audit checks split counts, unique recording identifiers, waveform validity,
and file hashes. No new split, trimming, or augmentation is applied.

\begin{table}[ht]
\centering\small
\begin{tabular}{lrrr}
\toprule
Partition & Normal & Defective & Total\\
\midrule
Training & 576 & 0 & 576\\
Validation & 122 & 1,610 & 1,732\\
Test & 121 & 1,611 & 1,732\\
\bottomrule
\end{tabular}
\caption{Manifest partitions, matching the published protocol. Test labels are
reserved for final evaluation.}\label{tab:data}
\end{table}

\section{Method and implementation choices}
\textbf{Features.}
The paper's selected FFT length is 16,384 at 192\,kHz, corresponding to
85.33\,ms. Scaling by sampling rate and rounding to an even integer gives
$N=1{,}366$ at 16\,kHz (85.375\,ms), with hop $H=683$ (42.6875\,ms).
We compute a one-sided Hann-window STFT without centering or normalization,
then $X_{ft}=\log\max(|\operatorname{STFT}(x)_{ft}|,10^{-8})$,
giving $F=684$ frequency bins. Duration is preserved approximately; bandwidth
is reduced to 8\,kHz and the frequency dimension changes.

\textbf{Architecture.}
Non-affine BatchNorm standardizes frequency channels. The encoder channel
sequence is $684\!\to\!1024\!\to\!1024\!\to\!1024\!\to\!1024\!\to\!48$;
five transposed convolutions reverse this sequence. All kernels have size
three, stride one, and no padding. The encoder reduces time by ten frames and
the decoder restores it. Activations are LeakyReLU with slope 0.01, except for
a scalar output PReLU initialized to 0.25. The model has 23,380,701 trainable
parameters. BatchNorm statistics update only during training; its standardized
output, detached from the computation graph, is the reconstruction target.

\textbf{Chunking, objective, and score.}
Inputs are divided into 192-frame chunks, with a final end-aligned chunk when
needed. Short chunks are reflect-padded. Padding is masked out of the loss and
score, although it influences convolutions and training BatchNorm. Waveforms
shorter than one FFT window are reflect-padded; incomplete trailing STFT
windows are discarded. For chunk $c$, let $Z_c$ be the standardized target,
$\widehat Z_c$ its reconstruction, and $m_{ct}$ indicate valid frames. We minimize
\begin{equation}
\mathcal L=
\frac{\sum_{c,t,f}m_{ct}(\widehat Z_{cft}-Z_{cft})^2}
{F\sum_{c,t}m_{ct}}.
\end{equation}
At inference, squared errors are averaged across frequency and, for frames
covered by multiple chunks, across their reconstructions. The recording score
is the mean over its original frames. Higher scores indicate anomalies.
Thus, overlap does not give the final frames extra weight in recording scores.

\section{Experimental protocol and observed results}
We train in float32 for 50 epochs with batch size 16, Adam
($\beta_1=0.9$, $\beta_2=0.999$, $\epsilon=10^{-8}$), and no weight decay.
A cosine one-cycle schedule peaks at $10^{-4}$ after a 30\% warmup, with
initial division factor 25 and final division factor 10,000; Adam momentum
is fixed. Training uses 2,880 chunks, shuffled with the run seed, and validation
uses 8,580 chunks. Runs use seeds 42, 43, and 44 on an NVIDIA GeForce RTX
5070 Ti with PyTorch 2.14.1+cu130. Deterministic PyTorch algorithms are enabled.

For each seed, the checkpoint with maximum pooled validation ROC-AUC is
retained; ties retain the earlier epoch. Seed 42 is the designated primary
result and is not replaced by the best seed. After all runs finish, checkpoint
hashes are frozen in \texttt{selection.json} before test evaluation is enabled.
The primary metric pools all defects versus normal recordings. Per-defect
ROC-AUC compares each defect with the same normal validation recordings;
their unweighted mean is a separate macro metric.

\begin{table}[ht]
\centering\small
\begin{tabular}{lrrrr}
\toprule
Experiment & Epochs done & Best epoch & Val. AUC & Test AUC\\
\midrule
Published, 192\,kHz & 50 & --- & 0.8451 & 0.8460\\
Replication, seed 42 & 50/50 & 47 & 0.8155 & 0.8178\\
Replication, seed 43 & 50/50 & 40 & 0.8155 & 0.8174\\
Replication, seed 44 & 50/50 & 49 & 0.8164 & 0.8188\\
Mean $\pm$ sample SD & --- & --- & $0.8158\pm0.0005$ & $0.8180\pm0.0007$\\
\bottomrule
\end{tabular}
\caption{Validation-selected results. Published values are from
Tables 7 and 9 of~\cite{stemmer2024}. All three replication runs and evaluations
are complete; test scores use the validation-selected, frozen checkpoints.
The mean summarizes individual runs, without ensembling.}\label{tab:results}
\end{table}

The primary test ROC-AUC is 0.0282 below the published test value;
its validation ROC-AUC is 0.0296 below the published validation value.
Across three seeds, the test sample standard deviation is 0.0007; this measures
variation from initialization and training order, rather than uncertainty from
sampling new recordings.
Seed 42's test macro defect AUC is 0.8497. Its highest per-defect test AUCs are
overlap (0.9993), crater cracks (0.9882), and undercut (0.9839); warping
(0.6162) and excessive penetration (0.6261) remain difficult.
We recomputed metrics from saved predictions, checked validation results against
run summaries and histories, and verified that test provenance and checkpoint
hashes match the frozen selection.

An independent real-data smoke check reduced MSE from 1.0001 to 0.2016 over
100 updates on a fixed batch from two normal training recordings and verified
checkpoint reload agreement. This checks optimization and serialization,
without measuring anomaly-detection generalization. The completed runs each
record approximately 7.3 minutes of run time; data preparation and the smoke
check are excluded.

\section{Limitations and reproducibility}
This experiment adapts a published architecture to a different input release.
Preserving STFT duration does not restore frequencies above 8\,kHz, and the
smaller frequency axis also changes parameter count. BatchNorm behavior,
padding, activation details, chunk length, and optimizer scheduling resolve
omissions in the source paper. Consequently, the validation gap cannot be
attributed to bandwidth alone. No ablation isolates these choices. Selection across
50 validation epochs can also favor validation-specific behavior. No confidence
interval, streaming latency measurement, or robustness
to a new welding station is established here. The source dataset also has
acknowledged label noise~\cite{stemmer2024}.

The implementation is in \path{src/intel_baseline/}; the full settings are in
\path{configs/intel_baseline.json}. After installing the environment and caching
the dataset archive, the repository-root command below audits the data, runs
or resumes training, freezes selection, evaluates the test partition, and
generates the report:
\begin{quote}\small\ttfamily bash scripts/run\_intel\_baseline.sh\end{quote}
Artifacts under \texttt{artifacts/intel\_baseline/} retain settings, histories,
best/latest checkpoints, predictions, and provenance. Random-generator,
optimizer, and scheduler states are saved for epoch-boundary resumption.
The dataset audit records revision \texttt{58e58e59beb7} and fingerprint prefix
\texttt{d250a3293b81d849}; full identifiers are in \path{data/intel_baseline/audit.json}.
Run metadata records code commit \texttt{0c2fba904be5} with a dirty working tree;
that commit alone therefore does not fully specify the executed source.
All three seed runs and frozen-checkpoint test evaluations are complete in this
artifact snapshot.

\begin{thebibliography}{1}
\small
\bibitem{stemmer2024}
Georg Stemmer, Jose A. Lopez, Juan A. Del Hoyo Ontiveros, Arvind Raju,
Tara Thimmanaik, and Sovan Biswas.
\newblock Unsupervised Welding Defect Detection Using Audio and Video.
\newblock \emph{arXiv preprint arXiv:2409.02290}, 2024.
\newblock \url{https://arxiv.org/abs/2409.02290}.
\end{thebibliography}
\end{document}
